\documentclass[11pt,letterpaper]{article}

\usepackage[margin=0.72in]{geometry}
\usepackage[T1]{fontenc}
\usepackage[utf8]{inputenc}
\usepackage{lmodern}
\usepackage{microtype}
\usepackage{enumitem}
\usepackage{tabularx}
\usepackage{titlesec}
\usepackage{xcolor}
\usepackage{hyperref}

\definecolor{linkblue}{HTML}{1F4E79}
\hypersetup{
  colorlinks=true,
  urlcolor=linkblue,
  linkcolor=linkblue
}
\urlstyle{same}

\setlength{\parindent}{0pt}
\setlength{\parskip}{0.18em}
\setlength{\tabcolsep}{0pt}
\pagestyle{empty}

\titleformat{\section}
  {\large\bfseries}
  {}
  {0pt}
  {}
  [\titlerule]
\titlespacing*{\section}{0pt}{0.95em}{0.45em}

\newcommand{\datedline}[2]{%
  \begin{tabularx}{\textwidth}{@{}X@{\hspace{0.8em}}r@{}}
    #1 & #2
  \end{tabularx}\par
}

\newcommand{\cvitem}[2]{%
  \textbf{#1:} #2\par
}

\newenvironment{cvitems}{%
  \begin{itemize}[leftmargin=1.2em,itemsep=0.03em,topsep=0.1em,parsep=0pt,partopsep=0pt]
}{%
  \end{itemize}
}

\newcommand{\paperentry}[4]{%
  \textbf{#1}\par
  #2\par
  \emph{#3}#4\par
}

\newcommand{\papersep}{\vspace{0.45em}}

\begin{document}

\begin{center}
  {\LARGE \textbf{Anxin (Bob) Guo}}\\[0.25em]
  \href{mailto:anxinbguo@gmail.com}{anxinbguo@gmail.com}
  \quad|\quad
  \href{https://anxin-bob-guo.org}{anxin-bob-guo.org}
\end{center}

\section{Education}

\datedline{\textbf{Northwestern University}, Evanston, Illinois}{September 2023 -- Present}
Ph.D. student in Computer Science.

\datedline{\textbf{Northwestern University}, Evanston, Illinois}{September 2019 -- June 2023}
B.A. in Mathematics and M.S. in Computer Science. GPA: 3.92/4.00.

\section{Research Interests}

Theoretical computer science and the theoretical foundations of machine learning, including learning theory, neural network expressivity, and information-theoretic interpretations of language models.

\section{Research Projects and Papers}

\paperentry
{AXIOM: Algorithmic Lemma Discovery via Formal Library Compression}
{With Ruize Xu, Yi Gu, Tianchong Jiang, Chenxiao Yang, and Zhiyuan Li.}
{Manuscript.}
{}
\papersep

\paperentry
{Copying Under Repetition: From Prefix Matching to Position Hashing}
{With Tianlong Huang, Chenxiao Yang, and Zhiyuan Li.}
{Manuscript.}
{}
\begin{cvitems}
  \item Formalized finite-horizon uniformity, an expressivity regime between non-uniform circuits and fully uniform transformers, where one model must handle all inputs up to length $N$ while resources may scale with $N$.
  \item Proved impossibility results for a class of transformers spanning NoPE and a practical implementation of ALiBi: they cannot correctly terminate the copy task even with depth $N^{0.01}$, near-logarithmic precision, and arbitrary width.
  \item Designed a modular attention head, a plug-in augmentation using value rotation with integer period $p$ to encode query-key distance through the value stream.
  \item Showed that the augmented architecture can copy arbitrary strings using logarithmic width and $\log \log N$ precision.
  \item Built diagnostic experiments showing that frontier 400B-parameter language models struggle to copy strings with extreme repetition.
\end{cvitems}
\papersep

\paperentry
{Hallucination is a Consequence of Space-Optimality: A Rate-Distortion Theorem for Membership Testing}
{With \href{https://ieor.columbia.edu/content/jingwei-li}{Jingwei Li}.}
{International Conference on Machine Learning (ICML), 2026, spotlight, top 2.2\%.}
{ \href{https://arxiv.org/abs/2602.00906}{arXiv}.}
\begin{cvitems}
  \item Defined membership testing as a unifying abstraction for approximate set membership, Bloom-filter-style error metrics, and factual errors in language models.
  \item Established the optimal space-error tradeoff in this framework, with an asymptotic lower bound governed by relative entropy, analogous to rate-distortion theory.
  \item Showed that optimally trained language models under cross-entropy loss with bounded memory may assign high confidence to unseen non-facts as a structural consequence of limited capacity.
  \item Extended known bounds for Bloom-type filters into a tight form, including regimes where false negatives are allowed.
\end{cvitems}
\papersep

\paperentry
{Agnostic Learning of Arbitrary ReLU Activation under Gaussian Marginals}
{With \href{https://users.cs.northwestern.edu/~aravindv/}{Aravindan Vijayaraghavan}.}
{Conference on Learning Theory (COLT), 2025.}
{ \href{https://arxiv.org/abs/2411.14349}{arXiv}; \href{https://proceedings.mlr.press/v291/guo25a.html}{conference version}; \href{https://youtu.be/bXVM10VRfcI?si=6dDID4kokQpQcaD2}{recorded talk}.}
\begin{cvitems}
  \item Gave the first polynomial-time constant-factor approximation agnostic learner for arbitrarily biased ReLU neurons under Gaussian marginals.
  \item Established a formal separation between Statistical Query (SQ) and Correlational Statistical Query (CSQ) models for this problem, proving limitations of gradient-based algorithms in this setting.
  \item Designed a method combining Thresholded PCA and Reweighted Gradient Descent to overcome limitations of standard gradient-based approaches.
\end{cvitems}
\papersep

\paperentry
{To Store or Not to Store: A Graph-Theoretic Approach for Dataset Versioning}
{With \href{https://ieor.columbia.edu/content/jingwei-li}{Jingwei Li}, \href{https://pattaras.github.io/}{Pattara Sukprasert}, \href{https://www.mccormick.northwestern.edu/research-faculty/directory/profiles/khuller-samir.html}{Samir Khuller}, \href{https://www.cs.umd.edu/~amol/}{Amol Deshpande}, and \href{https://research.adobe.com/person/koyel-mukherjee/}{Koyel Mukherjee}.}
{IEEE International Parallel and Distributed Processing Symposium (IPDPS), 2024.}
{ \href{https://arxiv.org/abs/2402.11741}{arXiv}; \href{https://ieeexplore.ieee.org/document/10579114}{conference version}.}
\begin{cvitems}
  \item Formulated and analyzed approximation algorithms and heuristics for optimizing the storage-retrieval cost tradeoff in large-scale dataset versioning systems.
  \item Designed provably near-optimal algorithms for tree-like graphs by exploiting low-treewidth structure.
  \item Implemented and benchmarked proposed heuristics on experiments based on real-world GitHub repository data, improving performance by up to 1000x over prior methods.
\end{cvitems}

\section{Awards and Honors}

\datedline{\textbf{Ph.D. Student Research Award}, Computer Science Department, Northwestern University}{2024 -- 2025}
\datedline{\textbf{Barris Award for Outstanding TA}, CS 335: Introduction to the Theory of Computing}{Fall 2024}
\datedline{\textbf{Junior Career Award in Mathematics}, Mathematics Department, Northwestern University}{2021 -- 2022}

\section{Research and Teaching}

\datedline{\textbf{IDEAL-funded summer research exchange}, University of Chicago}{Summer 2026}
Hosted by \href{https://www.stat.uchicago.edu/~chaogao/}{Chao Gao}.

\datedline{\textbf{Research Intern / IDEAL-funded summer research exchange}, TTIC}{Summer 2025}
Toyota Technological Institute at Chicago. Hosted by \href{https://zhiyuanli.ttic.edu/}{Zhiyuan Li}.
\begin{cvitems}
  \item Initiated the transformer copying project during the internship.
  \item Built experiments validating the failure mode of frontier language models and the proposed architecture.
  \item Presented at a reading group.
\end{cvitems}

\datedline{\textbf{Teaching Assistant}, CS 336: Design and Analysis of Algorithms}{Spring 2026}

\datedline{\textbf{Teaching Assistant}, CS 335: Introduction to the Theory of Computing}{Fall 2024}
\begin{cvitems}
  \item Gave a guest lecture on busy beaver numbers in addition to regular TA duties.
  \item Awarded the Barris Award for Outstanding TA.
\end{cvitems}

\section{Service and Programs}

\datedline{\textbf{Program Committee / Reviewer}, Reliable ML from Unreliable Data, NeurIPS Workshop}{2025}
\datedline{\textbf{Volunteer}, 65th IEEE Symposium on Foundations of Computer Science (FOCS)}{2024}
\datedline{\textbf{Ross Mathematics Program Asia}, student, junior counselor, counselor}{2018 -- 2020}
\datedline{\textbf{Directed Reading Program}, Linear Representations of Finite Groups}{Spring 2021}
Read with \href{https://wenyuanli1995-math.github.io/}{Wenyuan Li}.

\section{Skills}

\cvitem{Technology}{C++, Python (NumPy, PyTorch, pandas, matplotlib), Mathematica, and Gurobi Optimizer.}

\section{Writing}

I have written two Chinese-language articles on theoretical computer science for
\href{https://zhuanlan.zhihu.com/p/643661983}{my Zhihu column}.

\end{document}
